% TODO make hw stylesheet
\documentclass[10pt]{article}
\usepackage[letterpaper,top=1in]{geometry}
\usepackage[dvipsnames,svgnames]{xcolor}
\usepackage[shortlabels]{enumitem}
\usepackage[framemethod=TikZ]{mdframed}
\usepackage{amsmath,amssymb,amsthm}
\usepackage[colorlinks]{hyperref}
\usepackage{multicol}
\usepackage{tabularx}

\hypersetup{
	colorlinks=true,
	linkcolor=blue,
	filecolor=magenta,
	urlcolor=magenta,
	pdftitle={MAT 324 HW}
}

\usepackage{mathtools}
\usepackage{thmtools}
\usepackage[textsize=footnotesize]{todonotes}
\usepackage{titling}

\reversemarginpar
\mdfdefinestyle{mdbluebox}{roundcorner=10pt,innerbottommargin=9pt,
	linecolor=blue,backgroundcolor=TealBlue!5,}
\declaretheoremstyle[headfont=\sffamily\bfseries\color{MidnightBlue},
mdframed={style=mdbluebox},]{thmbluebox}

\mdfdefinestyle{mdbloobox}{frametitlefont=\bfseries,innerbottommargin=8pt,
	nobreak=true,backgroundcolor=TealBlue!5,linecolor=blue,}
\declaretheoremstyle[headfont=\bfseries\color{MidnightBlue},
mdframed={style=mdbloobox},headpunct={\\[3pt]},postheadspace=0pt,]{thmbloobox}

\mdfdefinestyle{mdredbox}{frametitlefont=\bfseries,innerbottommargin=8pt,
	nobreak=true,backgroundcolor=Salmon!5,linecolor=RawSienna,}
\declaretheoremstyle[headfont=\bfseries\color{RawSienna},
mdframed={style=mdredbox},headpunct={\\[3pt]},postheadspace=0pt,]{thmredbox}

\mdfdefinestyle{mdgreenbox}{linecolor=ForestGreen,backgroundcolor=ForestGreen!5,
	linewidth=2pt,rightline=false,leftline=true,topline=false,bottomline=false,}
\declaretheoremstyle[headfont=\bfseries\sffamily\color{ForestGreen!70!black},
mdframed={style=mdgreenbox},headpunct={ --- },]{thmgreenbox}

\mdfdefinestyle{mdorangebox}{linecolor=RedOrange,backgroundcolor=RedOrange!5,
	linewidth=2pt,rightline=false,leftline=true,topline=false,bottomline=false,}
\declaretheoremstyle[headfont=\bfseries\sffamily\color{RedOrange!70!black},
mdframed={style=mdorangebox},headpunct={ --- },]{thmorangebox}

\mdfdefinestyle{mdblackbox}{linecolor=black,backgroundcolor=RedViolet!5!gray!5,
	linewidth=3pt,rightline=false,leftline=true,topline=false,bottomline=false,}
\declaretheoremstyle[mdframed={style=mdblackbox}]{thmblackbox}

\declaretheoremstyle[spaceabove=6pt,spacebelow=6pt]{basehead}

\declaretheorem[style=basehead,name=Problem]{problem}
\declaretheorem[style=thmredbox,name=Problem,numbered=no]{problem*}
\declaretheorem[style=thmbloobox,name=Setup]{setup} % for config geo
\declaretheorem[style=thmredbox,name=Example,sibling=problem]{example}
\declaretheorem[style=thmredbox,name=$\bigstar$ Problem,sibling=problem]{niceprob}
\declaretheorem[style=thmbluebox,name=Theorem,sibling=problem]{theorem}
\declaretheorem[style=thmbluebox,name=Corollary,sibling=theorem]{corollary}
\declaretheorem[style=thmbluebox,name=Proposition,sibling=theorem]{prop}
\declaretheorem[style=thmbluebox,name=Lemma,numberwithin=problem]{lemma}
\declaretheorem[style=thmbluebox,name=Theorem,numbered=no]{theorem*}
\declaretheorem[style=thmbluebox,name=Lemma,numbered=no]{lemma*}
\declaretheorem[style=thmgreenbox,name=Claim,numberwithin=problem]{claim}
\declaretheorem[style=thmgreenbox,name=Claim,numbered=no]{claim*}
\declaretheorem[style=thmblackbox,name=Remark,numberwithin=problem]{remark}
\declaretheorem[style=thmblackbox,name=Remark,numbered=no]{remark*}
\declaretheorem[style=thmgreenbox,name=Definition,sibling=theorem]{definition}
\declaretheorem[style=thmgreenbox,name=Definition,numbered=no]{definition*}

\providecommand{\ol}{\overline}
\providecommand{\ceil}[1]{\left\lceil #1 \right\rceil}
\providecommand{\floor}[1]{\left\lfloor #1 \right\rfloor}
\newcommand{\dd}{\mathrm{d}}
\providecommand{\ul}{\underline}
\providecommand{\eps}{\varepsilon}
\providecommand{\half}{\frac{1}{2}}
\providecommand{\CC}{\mathbb C}
\providecommand{\FF}{\mathbb F}
\providecommand{\NN}{\mathbb N}
\providecommand{\QQ}{\mathbb Q}
\providecommand{\RR}{\mathbb R}
\providecommand{\ZZ}{\mathbb Z}
\providecommand{\dg}{^\circ}
\providecommand{\ii}{\item}
\providecommand{\setdiff}{\smallsetminus}
\providecommand{\alert}[1]{{\sffamily\textbf{\textcolor{blue}{#1}}}}
\providecommand{\inv}{^{-1}}
\providecommand{\mb}[1]{\mathbf{#1}}

\newenvironment{soln}{\begin{proof}[Solution]}{\end{proof}}

\providecommand{\pow}{\operatorname{Pow}}
\providecommand{\vocab}[1]{{\textbf{\textcolor{ForestGreen}{#1}}}}
\providecommand{\lcm}{\operatorname{lcm}}
\providecommand{\abs}[1]{\left|#1\right|}

\usepackage{mathrsfs}

% place title at top right
\setlength{\droptitle}{-8ex}
\pretitle{\begin{flushright}\Large\bfseries}
\posttitle{\par\end{flushright}}
\preauthor{\begin{flushright}}
\postauthor{\end{flushright}}
\predate{\begin{flushright}}
\postdate{\end{flushright}}

\title{MAT 324 HW02}
\author{\large Aditya Pahuja}
\date{\large Due 09/11}
\begin{document}
\maketitle

\begin{problem}
    Let $X$, $Y$ be sets, $\mathcal S\subseteq 2^X$ and $\mathcal T\subseteq 2^Y$
    be $\sigma$-algebras on $X$ and $Y$, respectively,
    and $f\colon X\to Y$ be a map.
    Suppose $\mathcal T$ is the $\sigma$-algebra generated by
    some subcollection $\mathcal T_0\subseteq\mathcal T$ and
    $f\inv(B)\in\mathcal S$ for every $B\in\mathcal T_0$.
    Show that $f\inv(B)\in\mathcal S$ for every $B\in\mathcal T$.
\end{problem}

\begin{soln}
    Let $\mathcal A = \{ A\subseteq Y : f\inv(A)\in\mathcal S\}$.
    Then $\mathcal A$ is a $\sigma$-algebra on $Y$ because
    (1) $f\inv(\varnothing) = \varnothing\in\mathcal S$,
    so $\varnothing\in \mathcal A$;
    (2) if $A\in\mathcal A$ then $f\inv(A)\in\mathcal S$
    hence $X\setminus f\inv(A) = f\inv(Y\setminus A)\in\mathcal S$,
    so $Y\setminus A\in\mathcal A$; and
    (3) if $A_1,A_2,\ldots\in\mathcal A$ then
    $f\inv(A_1),f\inv(A_2),\dots\in\mathcal S$ hence
    $\bigcup_{n=1}^\infty f\inv(A_n) =
    f\inv\left(\bigcup_{n=1}^\infty A_n\right)\in\mathcal S$,
    so $\bigcup_{n=1}^\infty A_n\in\mathcal A$.
    By assumption $\mathcal T_0\subseteq\mathcal A$,
    hence the $\sigma$-algebra generated by $\mathcal T_0$,
    namely $\mathcal T$, is also contained in $\mathcal A$,
    i.e. $f\inv(B)\in\mathcal S$ for each $B\in\mathcal T$.
\end{soln}

%\pagebreak

\begin{problem}
    Suppose $A\subseteq B\subseteq C$ are such that $A$, $C$ are Lebesgue
    and $\abs A = \abs C < \infty$. Show that $B$ is also Lebesgue.
\end{problem}

\begin{soln}
    By monotonicity $\abs A = \abs B = \abs C$.
    Since $\abs A$ and $\abs C$ are finite,
    $\abs{C\setminus A} = 0$, so $\abs{B\setminus A}\le\abs{C\setminus A}=0$
    implies that $B\setminus A$ is Lebesgue because it has outer measure zero.
    Then $B$ is Lebesgue because it is the union of Lebesgue sets
    $A$ and $B\setminus A$.
\end{soln}

%\pagebreak

\begin{problem}
    Let $A\subseteq\RR$ with $\abs A\ne 0$.
    Show that there exists $A'\subseteq A$ which is not Lebesgue measurable.
\end{problem}

\begin{soln}
    There exists an integer $N$ such that $\abs{A\cap[N-1,N+1]} > 0$;
    if such an $N$ didn't exist (so $\abs{A\cap [n-1,n+1]} = 0$ for all $n\in\ZZ$),
    then by subadditivity
    \begin{equation*}
	\abs A \le \sum_{n\in\ZZ} \abs{A\cap [n-1,n+1]} = 0.
    \end{equation*}
    Let $B = -N + (A\cap[N,N+1])\subseteq [-1,1]$, so $\abs B > 0$.

    The solution now proceeds nearly identically to Axler 2.18.
    Given numbers $x$ and $y$ in $B$, say that $x\sim y$ if $x-y\in\QQ$.
    This partitions $B$ into equivalence classes, and we write $[x]$
    for the equivalence class of $x\in B$.
    Define a set $S$ by taking one element
    from each of these equivalence class, i.e. $S\cap [x]$ has
    exactly one element for each $x\in B$.
    If $q_1$, $q_2$, \dots is an enumeration of $[-2,2]\cap\QQ$, we have
    \begin{equation*}
	B\subseteq \bigcup_{k=1}^\infty(q_k + S),
	\implies
	\abs B \le \sum_{k=1}^\infty \abs{q_k+S} = \sum_{k=1}^\infty \abs S,
    \end{equation*}
    implying $\abs S > 0$.
    On the other hand the translates $q_k + S$ are pairwise disjoint and
    \begin{equation*}
	\bigcup_{k=1}^\infty (q_k+S) \subseteq [-3,3].
    \end{equation*}
    Since the left-hand side is a countable union of pairwise disjoint sets,
    if $S$ were Lebesgue then we should be able to write
    \begin{equation*}
	\sum_{k=1}^\infty \abs S = \sum_{k=1}^\infty \abs{q_k+S}
	= \abs{\bigcup_{k=1}^\infty (q_k+S)} \le \abs{[-3,3]} = 6,
    \end{equation*}
    but this is a contradiction because $\abs S > 0$ implies
    the left-hand side is $\infty$, so $S$ is not Lebesgue.
\end{soln}

\pagebreak

\begin{problem}
    Let $A\subseteq\RR$ and $A^c\coloneq\RR\setminus A$
    be the complement of $A$. Show that $A$ is Lebesgue if and only if
    \begin{equation*}
	\abs E = \abs{E\cap A} + \abs{E \cap A^c} \text{ for all }E\subseteq\RR.
    \end{equation*}
\end{problem}

\begin{soln}
    Say that $A\subseteq\RR$ has the \emph{splitting property} if
    \begin{equation*}
	\abs E = \abs{E\cap A} + \abs{E \cap A^c} \text{ for all }E\subseteq\RR.
    \end{equation*}

    We first prove the \emph{only if} direction.
    Let $A$ be Lebesgue. It suffices to prove
    \begin{equation*}
	\abs E \ge \abs{E\cap A} + \abs{E\cap A^c}\text{ for all }E\subseteq\RR
    \end{equation*}
    because the reverse inequality holds by subadditivity.
    % If $\abs E = \infty$, then
    % \begin{equation*}
    %     \abs E = \abs{E\cap A} + \abs{E\cap A^c}
    % \end{equation*}
    % always holds because the right-hand side is at least as large as
    % $\abs E = \infty$, hence is $\infty$.
    % Assume now that $\abs E$ is finite.
    Since $\abs E$ is the infimum of $\abs U$ over
    all open sets $U$ (manifested as a union of
    countably many open intervals) that contain $E$,
    for a given $\eps>0$ we can find an open $U_\eps\subseteq\RR$
    such that $\abs{U_\eps}\le \abs E + \eps$
    This gives
    \begin{equation*}
	\abs E \ge \abs{U_\eps} - \eps
	= \abs{U_\eps\cap A} + \abs{U_\eps\cap A^c} - \eps
	\ge \abs{E\cap A} + \abs{E\cap A^c} - \eps
    \end{equation*}
    for all $\eps>0$, where
    $\abs{U_\eps} = \abs{U_\eps\cap A} + \abs{U_\eps\cap A^c}$
    because the sets $U_\eps\cap A$ and $U_\eps\cap A^c$
    are disjoint and Lebesgue. This implies that
    $A$ has the splitting property if it is Lebesgue.
    \\

    Suppose conversely that $A$ has the splitting property.
    If $\abs A <\infty$, then by definition of outer measure,
    given $\eps > 0$, we can find an open set $U_\eps$ for which
    $\abs{U_\eps} < \abs A + \eps$. Setting $E = U_\eps$ yields
    \begin{equation*}
	\abs{U_\eps} = \abs{U_\eps\cap A} + \abs{U_\eps\cap A^c}
	= \abs A + \abs{U_\eps\setminus A}
	\iff \abs{U_\eps\setminus A} = \abs{U_\eps} - \abs A < \eps.
    \end{equation*}
    The existence of such a $U_\eps$ is one of the criteria we proved in class
    for $A$ to be Lebesgue, sets of finite outer measure
    with the splitting property are Lebesgue.

    To finish, suppose that $A$ has the splitting property and $|A|=\infty$.
    We have proven above that $I_n \coloneq [n,n+1]$
    has the splitting property because it is Lebesgue, so
    \begin{align*}
	\abs{E\cap A\cap I_n} + \abs{E\cap (A\cap I_n)^c}
	&= \abs{E\cap A\cap I_n} + \abs{E\cap (A\cap I_n)^c\cap A} +
	\abs{E\cap (A\cap I_n)^c\cap A^c}\\
	&= \abs{E\cap A\cap I_n} + \abs{E\cap A\cap I_n^c} + \abs{E\cap A^c}\\
	&= \abs{E\cap A} + \abs{E\cap A^c} = \abs E,
    \end{align*}
    which proves that $A\cap I_n$ also has the splitting property.
    Since $A\cap I_n$ has finite outer measure and has the splitting property,
    it is Lebesgue, hence so is $A = \bigcup_{n\in\ZZ} (A\cap I_n)$.
\end{soln}

\pagebreak

\begin{problem}
    Let $X$ be a set, $\mathcal A\subseteq 2^X$, and
    $\ell\colon A\to[0,\infty]$ be a function such that
    $\varnothing\in\mathcal A$ and $\ell(\varnothing) = 0$.
    For each $A\subseteq X$, define
    \begin{equation*}
	\mathcal L_A \coloneq
	\left\{\sum_{n=1}^\infty \ell(I_n) : I_1,I_2,\ldots\in\mathcal A,\;
	A\subseteq\bigcup_{n=1}^\infty I_n\right\}\subseteq[0,\infty],\quad
	m(A) = \inf\mathcal L_A\in[0,\infty],
    \end{equation*}
    \begin{equation}\label{eq:cara-restriction}
	\mathcal M = \left\{A\subseteq X : m(E) = m(E\cap A) + m(E\cap A^c)
	\text{ for all }E\subseteq X\right\}.
    \end{equation}
    \begin{enumerate}[(a)]
	\ii Show that
	\begin{equation*}
	    m(\varnothing)=0;\; m(A)\le m(B)\text{ if }A\subseteq B;\;
	    m\left(\bigcup_{n=1}^\infty A_n\right)\le
	    \sum_{n=1}^\infty m(A_n)\text{ for all }A_1,A_2,\dots\subseteq X.
	    \stepcounter{equation}\tag{\theequation}\label{eq:om-props}
	\end{equation*}
	\ii Suppose in addition that $\mathcal A$ is an algebra on $X$ and
	\begin{equation*}
	    \ell\left(\bigcup_{n=1}^\infty I_n\right) = \sum_{n=1}^\infty\ell(I_n)
	\end{equation*}
	for all pairwise disjoint sequences of sets
	$I_1,I_2,\ldots\in\mathcal A$ such that
	$\bigcup_{n=1}^\infty I_n\in\mathcal A$.
	Show that $m|_{\mathcal A} = \ell$ and
	$\mathcal A\subseteq\mathcal M$.
\end{enumerate}
\end{problem}

\begin{soln}
    (a) $\mathcal L_{\varnothing}$ contains $0$, as we can cover $\varnothing$
    with the sequence of sets $I_1=I_2=\cdots=\varnothing$,
    so $\inf L_\varnothing = m(\varnothing) \le 0$,
    but of course $0$ is a lower bound for $\mathcal L_\varnothing$,
    so in fact $m(\varnothing) = 0$.

    Then, if $A\subseteq B$, a sequence of sets $I_1,I_2,\ldots\in\mathcal A$
    whose union contains $B$ must also contain $A$, meaning
    $\mathcal L_B\subseteq\mathcal L_A$; taking infima translates this to
    \begin{equation*}
	m(A) = \inf\mathcal L_A\le \inf\mathcal L_B = m(B).
    \end{equation*}

    Finally, let $A_1,A_2,\ldots\subseteq X$ be a sequence of subsets of $X$.
    For each $\eps > 0$ and $k\in\NN$, since $m(A_k) = \inf\mathcal L_{A_k}$,
    there exist sets $I^{(k)}_1,I^{(k)}_2,\ldots\in\mathcal A$
    such that $A_k\subseteq\bigcup_{k=1}^\infty I^{(k)}_n$ and
    \begin{equation*}
	\sum_{n=1}^\infty \ell(I^{(k)}_n)\le m(A_k) + \frac\eps{2^k}.
    \end{equation*}
    Summing over all $k$, we get
    \begin{equation*}
	\eps+\sum_{k=1}^\infty m(A_k) =
	\sum_{k=1}^\infty \left(m(A_k) + \frac\eps{2^k}\right)
	\ge \sum_{k=1}^\infty \sum_{n=1}^\infty \ell(I^{(k)}_n)
	\ge m\left(\bigcup_{k=1}^\infty A_k\right)
    \end{equation*}
    where the last inequality is a result of the fact that
    the intervals $\{I^{(k)}_n\}_{k,n\in\NN}$ cover $\bigcup_{k=1}^\infty A_k$.

    (b) Suppose that $A$ is in $\mathcal A$ and let $\eps>0$.
    Then, there exists a sequence $I_1,I_2,\ldots\in\mathcal A$
    such that $A\subseteq \bigcup_{n=1}^\infty I_n$ and
    \begin{equation*}
	\sum_{n=1}^\infty \ell(I_n) \le m(A) + \eps.
    \end{equation*}
    Let $I'_n = I_n\cap A\in\mathcal A$ and define $J_1 = I'_1$
    and $J_n = I'_n\setminus\bigcup_{k=1}^{n-1}I'_n\in\mathcal A$ for $n\ge 2$.
    The $J_n$ are pairwise disjoint and
    $\bigcup_{n=1}^\infty J_n = A\in\mathcal A$, so
    \begin{equation*}
	\ell(A) = \sum_{n=1}^\infty \ell(J_n)
	\le \sum_{n=1}^\infty \ell(I'_n)
	\le \sum_{n=1}^\infty \ell(I_n)
	\le m(A) + \eps.
    \end{equation*}
    This holds for all $\eps>0$, implying $\ell(A) \le m(A)$.
    On the other hand, $A$ can be covered by
    $A$ and countably many copies of $\varnothing$, implying that
    \begin{equation*}
        m(A) \le \ell(A) + \ell(\varnothing) + \ell(\varnothing) + \cdots
	= \ell(A).
    \end{equation*}
    Therefore, $m(A) = \ell(A)$ for all $A\in\mathcal A$.
    \\

    Suppose now that $E\subseteq X$ is arbitrary and $A\in\mathcal A$.
    Then, given $\eps>0$, there exist sets $I_1,I_2,\dots\in\mathcal A$
    whose union contains $A$ such that
    \begin{equation*}
	\sum_{n=1}^\infty \ell(I_n) \le m(E) + \eps.
    \end{equation*}
    Noting that $A^c\in\mathcal A$
    and so $I_n\cap A$, $I_n\cap A^c\in\mathcal A$,
    we thus get the inequality
    \begin{align*}
	m(E\cap A) + m(E\cap A^c)
	&\le m\left(\left(\bigcup_{n=1}^\infty I_n\right) \cap A\right)
	+ m\left(\left(\bigcup_{n=1}^\infty I_n\right)\cap A^c\right)\\
	&= m\left(\bigcup_{n=1}^\infty (I_n\cap A)\right)
	+ m\left(\bigcup_{n=1}^\infty (I_n\cap A^c)\right)\\
	&\le \sum_{n=1}^\infty \ell(I_n\cap A) + \ell(I_n\cap A^c) =
	\sum_{n=1}^\infty \ell(I_n) \le m(E) + \eps
    \end{align*}
    where the equality in the last line holds
    by additivity on the disjoint sets
    $I_n\cap A$, $I_n\cap A^c$
    (and we can pad the finite union into an infinite one
    with infinitely many copies of $\varnothing$).
    The resulting inequality holds
    for all $\eps>0$, so $m(E\cap A) + m(E\cap A^c) \le m(E)$.
    The reverse inequality is immediate by subadditivity,
    hence $m(E) = m(E\cap A) + m(E\cap A^c)$ for all $E\subseteq X$,
    i.e. $A\in\mathcal M$.
\end{soln}

\pagebreak

\begin{problem}
    Suppose $X$ is a set, $m\colon 2^X\to[0,\infty]$ is a function satisfying
    (2), and $\mathcal M\subseteq 2^X$ is as in (1).
    For $A\subseteq X$, let $A^c\coloneq X\setminus A$
    be the complement of $A$ in $X$. Show that
    \begin{enumerate}[(a)]
	\ii $m(A\cup B)\le m(A) + m(B)$ for all $A,B\subseteq X$;
	\ii $X\in\mathcal M$, $A^c\in\mathcal M$ if $A\in\mathcal M$,
	and $A\in\mathcal M$ if $m(A) = 0$;
	\ii if $A,B\in\mathcal M$ and $E\subseteq X$, then
	\begin{align*}
	    m(E) &= m(E\cap A\cap B) + m(E\cap A\cap B^c)
	    + m(E\cap A^c\cap B) + m(E\cap A^c\cap B^c)\\
	    &\ge m(E\cap(A\cup B)) + m(E\cap(A\cup B)^c);
	\end{align*}
	\ii $\mathcal M$ is an algebra on $X$ (not $\sigma$-algebra yet) and
	\begin{equation*}
	    m\left(\bigcup_{k=1}^n A_k\right) = \sum_{k=1}^n m(A_k)
	\end{equation*}
	for all sequences of pairwise disjoint sets
	$A_1,\dots, A_n\in\mathcal M$;
	\ii if $A_1,A_2,\ldots\in\mathcal M$ are pairwise disjoint, then
	\begin{equation*}
	    m\left(E\cap \bigcup_{n=1}^\infty A_n\right) =
	    \sum_{n=1}^\infty m(E\cap A_n)
	    \text{ for all $E\subseteq X$,}\quad
	    m\left(\bigcup_{n=1}^\infty A_n\right) = \sum_{n=1}^\infty m(A_n);
	\end{equation*}
	\ii $(X,\mathcal M,\mu\equiv m|_{\mathcal M})$
	is a complete measure space.
\end{enumerate}
\end{problem}

\begin{soln}
    (a) Take $A_1=A$, $A_2=B$, and, for integers $n>2$, $A_n = \varnothing$.
    Then,
    \begin{equation*}
	m(A\cup B) = m\left(A\cup B\cup \bigcup_{n=3}^\infty \varnothing\right)
	\le m(A) + m(B) + \sum_{n=3}^\infty m(\varnothing)
	= m(A) + m(B).
    \end{equation*}

    (b) $X\in\mathcal M$ because for any $E\subseteq X$,
    \begin{equation*}
	m(E\cap X) + m(E\cap X^c)
	= m(E) + m(\varnothing) = m(E).
    \end{equation*}
    If $A\in\mathcal M$, then $A^c\in\mathcal M$ because
    for any $E\subseteq X$,
    \begin{equation*}
	m(E\cap A^c) + m(E\cap(A^c)^c) = m(E\cap A^c) + m(E\cap A) = m(E).
    \end{equation*}
    If $A\in\mathcal M$ has $m(A) = 0$, then
    \begin{equation*}
        m(E)\le m(E\cap A) + m(E\cap A^c) = m(E\cap A^c) \le m(E)
    \end{equation*}
    where the first inequality is by (a)
    and the equality is by $m(E\cap A)\le m(A)$ implying $m(E\cap A) = 0$.\\

    (c) If $A,B\in\mathcal M$, then
    \begin{align*}
	m((E\cap A)\cap B) + m((E\cap A)\cap B^c) &= m(E\cap A)\\
	m((E\cap A^c)\cap B) + m((E\cap A^c)\cap B^c) &= m(E\cap A^c),
    \end{align*}
    so the left-hand sides sum to
    $m(E\cap A) + m(E\cap A^c) = m(E)$, proving the equality.
    The inequality follows from
    $A\cup B = (A\cap B)\cup(A\cap B^c)\cup(A^c\cap B)$,
    since subadditivty says that
    \begin{equation*}
	m(E\cap A\cap B) + m(E\cap A\cap B^c) + m(E\cap A^c\cap B^c)
	\ge m(E\cap(A\cup B))
    \end{equation*}
    and then we can add $m(E\cap(A^c\cap B^c)) = m(E\cap(A\cup B)^c)$
    to both sides.
    \pagebreak

    (d) $\mathcal M$ contains $X$ and is closed under complements by (b).
    By (c), for any $E\subseteq X$,
    \begin{equation*}
	m(E) \ge m(E\cap(A\cup B)) + m(E\cap(A\cup B)^c),
    \end{equation*}
    and the reverse inequality holds by (a), so
    $m(E) = m(E\cap(A\cup B)) + m(E\cap(A\cup B)^c)$;
    This is what it means for $A\cup B$ to be in $\mathcal M$.
    Suppose now that unions of up to $n-1$ sets in $\mathcal M$
    are also in $\mathcal M$; then, if $A_1,\dots,A_n\in\mathcal M$,
    so is 
    \begin{equation*}
	\bigcup_{k=1}^{n-1} A_k \cup A_n,
    \end{equation*}
    so by induction $\mathcal M$ is closed under arbitrary finite unions.

    Now, if $A,B\in\mathcal M$ are disjoint sets, then
    \begin{equation*}
	m(A\cup B) = m((A\cup B)\cap A) + m((A\cup B)\cap A^c)
	= m(A) + m(B),
    \end{equation*}
    i.e. additivity holds for two disjoint sets.
    If we suppose that additivity holds for $n-1$ pairwise disjoint
    sets in $\mathcal M$, then, given pairwise disjoint sets
    $A_1,\dots,A_n\in\mathcal M$,
    \begin{equation*}
	m\left(\bigcup_{k=1}^n A_k\right)
	= m\left(\bigcup_{k=1}^{n-1} A_k\right) + m(A_n)
	= \sum_{k=1}^{n-1}m(A_k) + m(A_n)
	= \sum_{k=1}^n m(A_n).
    \end{equation*}
    By induction, the measure of the union of any finite number of
    pairwise disjoint sets in $\mathcal M$
    is equal to the sum of the measures of the component sets.
    \\

    (e) Suppose $A_1,A_2,\ldots\in\mathcal M$ are pairwise disjoint.
    Then, for any positive integer $N$,
    \begin{equation*}
	m\left(E\cap\bigcup_{n=1}^\infty A_n\right) \ge
	m\left(E\cap\bigcup_{n=1}^N A_n\right).
    \end{equation*}
    I now claim that the right-hand side can be written as
    $\sum_{n=1}^N m(E\cap A_n)$.
    This follows from an induction argument:
    the version for two disjoint sets is shown by
    \begin{align*}
	m(E\cap(A\cup B)) &=
	m(E\cap(A\cup B)\cap A) + m(E\cap(A\cup B)\cap A^c)\\
			  &= m(E\cap A) + m(E\cap B)
    \end{align*}
    where $(A\cup B)\cap A^c = B$ because $A$ and $B$ are disjoint,
    and, assuming the result for unions of $N-1$ pairwise disjoint sets,
    we get that
    \begin{align*}
	m\left(E\cap\bigcup_{n=1}^N A_n\right)
	&= m\left(E\cap\bigcup_{n=1}^{N-1} A_n\right) + m(E\cap A_N)\\
	&= \sum_{n=1}^{N-1} (m(E\cap A_n)) + m(E\cap A_N)
	= \sum_{n=1}^N m(E\cap A_n)
    \end{align*}
    In summary, we have shown that
    \begin{equation*}
	m\left(E\cap\bigcup_{n=1}^\infty A_n\right) \ge
	\sum_{n=1}^N m(E\cap A_n).
    \end{equation*}
    In the limit as $N$ grows arbitrarily large,
    \begin{equation*}
	m\left(E\cap\bigcup_{n=1}^\infty A_n\right) \ge
	\lim_{N\to\infty}\sum_{n=1}^N m(E\cap A_n)
	= \sum_{n=1}^\infty m(E\cap A_n).
    \end{equation*}
    The reverse inequality holds by countable subadditivity,
    so in fact the inequality is sharp.
    Then, we take $E = \bigcup_{n=1}^\infty A_n$ to get
    \begin{equation*}
	m\left(\bigcup_{n=1}^\infty A_n\right) =
	\sum_{n=1}^\infty m\left(\bigcup_{k=1}^\infty A_k\cap A_n\right) =
	\sum_{n=1}^\infty m(A_n).
    \end{equation*}

    (f) Suppose $A_1,A_2,\ldots\in\mathcal M$.
    If $m(E\cap\bigcap_{n=1}^\infty A_n) = \infty$, then
    by the assumption (2) we have $m(E) = \infty$ as well,
    in which case
    $m(E\cap\bigcup_{n=1}^\infty A_n) + m(E\cap\bigcap_{n=1}^\infty A_n^c) = m(E)$.
    Assume now that $m(E\cap\bigcap_{n=1}^\infty A_n) < \infty$.
    Define $B_1 = A_1$ and $B_n = A_n\setminus\bigcup_{k=1}^{n-1}A_k$ for $n>1$,
    and note that since $\mathcal M$ is an algebra,
    $\bigcup_{k=1}^n A_n = \bigcup_{k=1}^n B_n\in\mathcal M$ for all $n\in\NN$.
    By (e), since the $B_n$ are pairwise disjoint,
    \begin{align*}
	m\left(E\cap\bigcup_{n=1}^\infty A_n\right) +
	m\left(E\cap\bigcap_{n=1}^\infty A_n^c\right)
	&= m\left(E\cap\bigcup_{n=1}^\infty B_n\right) +
	m\left(E\cap\bigcap_{n=1}^\infty B_n^c\right)\\
	&= \sum_{n=1}^\infty (m(E\cap B_n)) +
	m\left(E\cap\bigcap_{n=1}^\infty B_n^c\right)
    \end{align*}
    The sum is finite, so given $\eps>0$
    there exists $N$ such that
    $\sum_{n=1}^\infty m(E\cap B_n) < \eps+\sum_{n=1}^N m(E\cap B_n)$.
    Thus
    \begin{align*}
	\sum_{n=1}^\infty (m(E\cap B_n)) +
	m\left(E\cap\bigcap_{n=1}^\infty B_n^c\right)
	&\le \sum_{n=1}^N (m(E\cap B_n)) +
	m\left(E\cap\bigcap_{n=1}^N B_n^c\right) + \eps\\
	&= m\left(E\cap\bigcup_{n=1}^N B_n\right) +
	m\left(E\cap\bigcap_{n=1}^N B_n^c\right) + \eps = m(E) + \eps
    \end{align*}
    where the first equality follows from part (e),
    since we can pad the union of the (pairwise disjoint) sets
    $B_1$, \dots, $B_N$ into a countable union with copies of $\varnothing$;
    and the second equality follows from $\bigcup_{k=1}^n B_n$
    being in $\mathcal M$ as stated earlier. In summary,
    \begin{equation*}
	m\left(E\cap\bigcup_{n=1}^\infty A_n\right) +
	m\left(E\cap\bigcap_{n=1}^\infty A_n^c\right)
	\le m(E) + \eps
    \end{equation*}
    for every $\eps>0$, so the left-hand side is bounded above by $m(E)$.
    Yet again the reverse inequality is true by subadditivity, so
    $m(E\cap\bigcup_{n=1}^\infty A_n) + m(E\cap\bigcap_{n=1}^\infty A_n^c)
    = m(E)$ for all $E$, which is what it means for
    the union $\bigcup_{n=1}^\infty A_n$ to be in $\mathcal M$;
    this proves that $\mathcal M$ is a $\sigma$-algebra.

    To show that $\mu \equiv m|_{\mathcal M}$ is a measure,
    we have $\mu(\varnothing) = 0$ by assumption and
    if $A_1,A_2,\ldots\in\mathcal M$ are pairwise disjoint, then by (e)
    \begin{equation*}
	\mu\left(\bigcup_{n=1}^\infty A_n\right) =
	\sum_{n=1}^\infty \mu(A_n).
    \end{equation*}
    This completes the proof that $(X,\mathcal M, \mu)$ is a measure space.\\

    Finally, this measure space is complete:
    given $A\in\mathcal M$ with $\mu(A) = 0$, any subset
    $B\subseteq A$ has $m(B)\le m(A) = \mu(A)$,
    which implies $m(B) = 0$, so by part (b), $B\in\mathcal M$.
\end{soln}










\end{document}
